\HMTNarrativeHeading{Une unité, deux réalisations}

La double projection organise cet article dès sa première construction. Un
même état conserve les données qui permettent de préciser une valeur et les données
qui permettent de déterminer une incidence. La première réalisation réunit
des cylindres compatibles jusqu’à individualiser sa lecture numérique ; la seconde
transporte les relations de frontière. La thèse étudie leur provenance
commune et les opérations qui maintiennent leur correspondance pendant le
raffinement. Nombre et configuration géométrique reçoivent ainsi une généalogie
explicite.

L’unité se constitue au moyen de positions, d’opérations et de relations de
transport. APP dispose les deux directions nonadiques sur un même
support et évalue ses feuillets additif et multiplicatif. Chaque évaluation
conserve le résidu et son quotient. TRIT distingue l’orientation et le
régime local ; le TPK compose ces données avec le déplacement, la retenue,
la frontière et le prolongement. La structure discrète conjointe du
continu réside dans cette construction enrichie. Ses réalisations
ultérieures reçoivent des opérations déjà déterminées.

La connexion nonadique possède ici une fonction constitutive. Un tour compose les transitions de l'état complet. La phase observable revient et le registre incorpore le parcours accompli. Cette différence permet de reconnaître une régularité tandis que l'histoire gagne en profondeur. La continuité du raffinement conserve les préfixes qui rendent les deux réalisations compatibles. Chaque nouvelle précision appartient à la même histoire dont proviennent les précédentes.

Le sens de la conservation évolutive apparaît à ce point. Une quantité
normalisée peut demeurer égale tandis que changent la distribution, la
résolution ou la représentation. Récupérer l’état exige d’identifier quelles
données accompagnent ce changement. L’arithmétique conserve les quotients ; le transport
conserve l’orientation et les tours ; l’incidence conserve les relations ;
l’archive opératorielle conserve les compléments nécessaires pour inverser la
lecture. L’article développe les applications qui articulent ces fonctions.

L'extrême et moyenne raison holographique exprime cette question d'unité, de proportion et de transformation. Sa formulation inclut la complétion dimensionnelle et le raffinement réversible, et atteint ensuite les bilans d'états et d'observables. La constante de couplage occupe une position déterminée dans ce parcours : le registre engendré relie la lecture arithmétique, l'incidence et la récupération opératorielle. Les démonstrations suivantes construisent les antécédents de chacune de ces fonctions.

